\subsection{Homotopies}

DEFINITION 4. --- Let (C, d) and (C', d') be two complexes, and let f and g be two morphisms from C to C'. A homotopy connecting f to g is any graded A-homomorphism s of degree 1 from C to C' such that \(g - f = d' \circ s + s \circ d\) .

We say that f and g are homotopic if there exists a homotopy connecting f to g.

If \(h\) is a third morphism of \(C\) to \(C'\) and if \(s\) (resp. \(t\) ) is a homotopy connecting \(f\) to \(g\) (resp. \(g\) to \(h\) ), then \(s + t\) is a homotopy connecting \(f\) to \(h\) ; consequently, the relation “f and g are two homotopic morphisms from C to C'” is an equivalence relation, whose classes are called the homotopy classes of morphisms from C to C'.

Given two topological spaces X and Y, and a continuous map \(f: \mathbf{X} \to \mathbf{Y}\) , we define a linear map \(f_{*}\) of the singular complex (cf.~no. 3) C(X, A) into C(Y, A) by setting \(f_{*}(e_{s}) = e_{f \circ s}\) for \(s \in \Sigma_{n}(\mathbf{X})\) . This map is a morphism of complexes.

Two continuous maps \(f\) and \(g\) from X to Y are said to be topologically homotopic if there exists a continuous map \(h\) of \([0, 1] \times X\) into Y such that \(h(0, x) = f(x)\) and \(h(1, x) = g(x)\) for all \(x \in X\) . We show that, if \(f\) and \(g\) are topologically homotopic, the morphisms \(f_*\) and \(g_*\) are homotopic in the sense of Definition 4 above. It is this fact that is the origin of the terminology used in algebra. *

PROPOSITION 3. --- If f and g are two homotopic morphisms from C to C', then H(f) = H(g).

Let s be a homotopy connecting f to g. We have

\[
(g - f) (\mathbf {Z} (\mathbf {C})) = (d ^ {\prime} \circ s + s \circ d) (\mathbf {Z} (\mathbf {C})) = (d ^ {\prime} \circ s) (\mathbf {Z} (\mathbf {C})) \subset \mathrm{B} (\mathbf {C} ^ {\prime}),
\]

hence \(\mathrm{H}(g - f) = 0\) and \(\mathrm{H}(g) = \mathrm{H}(f)\) .

COROLLARY. --- A morphism homotopic to a homologism is a homologism.

PROPOSITION 4. --- Let C, C', D, D' be four complexes, \(f: \mathbf{C} \to \mathbf{C}'\) , \(g: \mathbf{C} \to \mathbf{C}'\) , \(u: \mathbf{D} \to \mathbf{C}\) , \(v: \mathbf{C}' \to \mathbf{D}'\) four morphisms. If s is a homotopy relating \(f\) to \(g\) , then \(v \circ s \circ u\) is a homotopy connecting \(v \circ f \circ u\) to \(v \circ g \circ u\) . If f and g are homotopic, then so are \(v \circ f \circ u\) and \(v \circ g \circ u\) .

That is clear.

COROLLARY. --- Let C, C′, C″ be three complexes, and f and g two morphisms from C to C′, \(f_{1}\) and \(g_{1}\) two morphisms from C' to C''. If s and \(s_{1}\) are homotopies connecting f to g and \(f_{1}\) to \(g_{1}\) respectively, then \(s_{1} \circ f + g_{1} \circ s\) is a homotopy connecting \(f_{1} \circ f\) to \(g_{1} \circ g\) . If f and \(f_{1}\) are homotopic to g and \(g_{1}\) respectively, then \(f_{1} \circ f\) is homotopic to \(g_{1} \circ g\) .

Indeed, \(s_{1} \circ f\) connects \(f_{1} \circ f\) to \(g_{1} \circ f\) and \(g_{1} \circ s\) connects \(g_{1} \circ f\) to \(g_{1} \circ g\) .

DEFINITION 5. --- A morphism of complexes \(f: \mathbb{C} \to \mathbb{C}'\) is called a homotopy equivalence if there exists a morphism \(f': \mathbb{C}' \to \mathbb{C}\) such that \(f' \circ f\) and \(f \circ f'\) are homotopic to \(1_{\mathbb{C}}\) and \(1_{\mathbb{C}'}\) respectively.

It is clear that \(f'\) is then also a homotopy equivalence; one also says that \(f'\) is the inverse of \(f\) up to homotopy. If \(f'\) and \(f_1'\) are both inverses of \(f\) up to homotopy, then \(f'\) and \(f_1'\) are homotopic (indeed, by the preceding corollary, \(f_1' = f_1' \circ 1_{\mathbf{C}'}\) is homotopic to \(f_1' \circ f \circ f'\) , hence to \(1_{\mathbf{C}} \circ f' = f'\) ).

PROPOSITION 5. — A homotopism is a homologism; a morphism composed of homotopisms is a homotopism. A morphism homotopic to a homotopism is a homotopism.

Let \(f: \mathbf{C} \to \mathbf{C}'\) and \(f_1: \mathbf{C}' \to \mathbf{C}''\) be homotopisms of complexes, \(f': \mathbf{C}' \to \mathbf{C}\) and \(f_1': \mathbf{C}'' \to \mathbf{C}'\) morphisms inverse up to homotopy. We have

\[
\mathrm{H} (f ^ {\prime}) \circ \mathrm{H} (f) = \mathrm{H} (f ^ {\prime} \circ f) = \mathrm{H} (1 _ {\mathrm{C}}) = 1 _ {\mathrm{H(C)}} \quad (\text { prop. } 3)
\]

and likewise \(\mathrm{H}(f) \circ \mathrm{H}(f') = 1_{\mathrm{H}(\mathbb{C}')}\) , hence \(\mathrm{H}(f)\) is bijective and \(f\) is a homologism. On the other hand, \((f' \circ f_1') \circ (f_1 \circ f)\) is homotopic to \(f' \circ 1_{\mathbb{C}'} \circ f\) (prop. 4), hence to \(1_{\mathbb{C}}\) ; similarly, \((f_1 \circ f) \circ (f' \circ f_1')\) is homotopic to \(1_{\mathbb{C}''}\) and \(f_1 \circ f\) is a homotopism. Finally, if \(g: \mathbb{C} \to \mathbb{C}'\) is a morphism homotopic to \(f\) , \(f' \circ g\) is homotopic to \(f' \circ f\) , hence to \(1_{\mathbb{C}}\) ; similarly, \(g \circ f'\) is homotopic to \(f \circ f'\) , hence to \(1_{\mathbb{C}'}\) and \(g\) is a homotopism.

COROLLARY. --- Let C, C', D, D' be four complexes, \(f: \mathbf{C} \to \mathbf{C}'\) a morphism, \(u: \mathbf{D} \to \mathbf{C}\) and \(v: \mathbf{C}' \to \mathbf{D}'\) homotopisms. In order that \(v \circ f \circ u\) be a homotopism (resp. a homologism), it is necessary and sufficient that \(f\) be one.

If \(f\) is a homotopism (resp. a homologism), then \(v \circ f \circ u\) is composed of homotopisms (resp. homologisms), hence is one. Conversely, let \(\overline{u}\) and \(\overline{v}\) be morphisms inverse to \(u\) and \(v\) up to homotopy; then \(\overline{v} \circ (v \circ f \circ u) \circ \overline{u}\) is homotopic to \(f\) by prop. 4; whence the conclusion by prop. 5, and the corollary of prop. 3.

We say that the complex C is homotopic to zero if \(1_{C}\) is homotopic to the zero map, that is, if there exists a graded A-endomorphism s of degree 1 of C such that \(1_{C} = s \circ d + d \circ s\) . This also amounts to saying that the unique morphism \(0 \to C\) (resp. \(C \to 0\) ) is a homotopism. A complex homotopic to zero has zero homology (prop. 5).

Example. --- Let \(u: \mathbf{M} \to \mathbf{N}\) and \(v: \mathbf{N} \to \mathbf{P}\) be homomorphisms of A-modules such that \(v \circ u = 0\) ; let C be the complex such that \(\mathbf{C}_2 = \mathbf{M}\) , \(\mathbf{C}_1 = \mathbf{N}\) , \(\mathbf{C}_0 = \mathbf{P}\) , \(\mathbf{C}_i = 0\) for \(i \neq 0, 1, 2, d_2 = u, d_1 = v, d_i = 0\) for \(i \neq 1, 2\) . Then C has zero homology if and only if the sequence \(0 \to \mathbf{M} \xrightarrow{u} \mathbf{N} \xrightarrow{v} \mathbf{P} \to 0\) is exact. It is homotopic to zero if and only if this sequence is split. Indeed, saying that C is homotopic to zero means that there exist A-homomorphisms \(s: \mathbf{P} \to \mathbf{N}\) and \(t: \mathbf{N} \to \mathbf{M}\) such that \(v \circ s = 1_{\mathbf{P}}, s \circ v + u \circ t = 1_{\mathbf{N}}, t \circ u = 1_{\mathbf{M}}\) ; this implies that the sequence is split; conversely, if \(s\) is an A-linear section of \(v\) , we define \(t\) by \(u \circ t = 1_{\mathbf{N}} - s \circ v\) , which is possible since \(v \circ (1_{\mathbf{N}} - s \circ v) = v - v \circ s \circ v = 0\) .

